\documentclass[11pt]{article}
\usepackage[margin=1in]{geometry}
\usepackage{booktabs}
\usepackage{longtable}
\usepackage{hyperref}
\usepackage{xcolor}
\usepackage{graphicx}
\usepackage{textcomp}
\hypersetup{colorlinks=true,linkcolor=blue,urlcolor=blue,citecolor=blue}

\title{Replication Report --- Subedi et al.\ (2019),\\
\textit{P.\ aeruginosa} PA34 Accessory Genome}
\author{Ollie (OpenClaw AI) --- X-100 Replication Project, BVBRC set (index 92)}
\date{2026-07-04}

\begin{document}
\maketitle

\section*{Header}
\textbf{Paper:} Subedi D, Kohli GS, Vijay AK, Willcox M, Rice SA.
\textit{Accessory genome of the multi-drug resistant ocular isolate of Pseudomonas aeruginosa PA34.}
\textbf{PLoS ONE} 14(4):e0215038 (April 15, 2019). \\
\textbf{DOI:} \href{https://doi.org/10.1371/journal.pone.0215038}{10.1371/journal.pone.0215038} \textperiodcentered{}
\textbf{PMID:} 30986237 \textperiodcentered{} \textbf{PMC:} PMC6464166 \textperiodcentered{}
Open access (CC BY 4.0 / PLOS). \\
\textbf{Deposited data:} GenBank CP032552 (chromosome), MH547560 (pMKPA34-1), MH547561 (pMKPA34-2);
BioProject PRJNA431326; BioSample SAMN08435059; BV-BRC \texttt{genome\_id} 287.6355. \\
\textbf{Verdict:} \textbf{PARTIAL REPLICATION (strong).}

\section{Paper summary}
The authors sequence \textit{P.\ aeruginosa} strain PA34---a multi-drug-resistant microbial-keratitis
isolate collected in 1997 in Hyderabad, India---by Illumina + Oxford Nanopore MinION hybrid assembly
with PCR gap closure, producing a closed 6.81 Mbp chromosome (66.1\% GC, 6{,}462 CDS) plus
\textbf{pMKPA34-1} (95.4 kbp, 57\% GC, 98 CDS) and \textbf{pMKPA34-2} (26.8 kbp, 61\% GC, 33 CDS).
A Roary pan-genome comparison against PAO1, PA14, and VRFPA04 identifies 7{,}643 orthologs
(5{,}078 core; PA34 accessory 1{,}213 with 543 unique). MAUVE + comparative genomics identifies
24 regions of genomic plasticity, including two novel islands \textbf{MKPA34-GI1} (68.6 kbp,
chromate + mercury resistance) and \textbf{MKPA34-GI2} (35.9 kbp, phage MP38), three ICEs
(RGP5, RGP29, RGP41 = pKLC102/PAPI-1-like), a functional \textbf{exoU} cytotoxin island (RGP7),
an \textbf{AAC(3)-IId} aminoglycoside-resistance gene in a 125 kbp island (RGP23 --- also carrying
tunicamycin + copper resistance + phage gp37), and \textbf{two independent mercury-resistance operons}.
The plasmids together carry six AMR genes (dfrA15, cmlA1, APH(3\textquotedbl)-Ib, APH(6)-Id,
bla$_{\text{NPS-1}}$, acrB on pMKPA34-1 in a Tn3-like transposon with class-I integron
\textbf{In1427}; mepA + full Tn7 module on pMKPA34-2). PA34 is high-cytotoxic (Fig 5) and Hg-tolerant
(Fig 6, $p<0.05$ vs PAO1). CRISPR-Cas absent.

\section{Claims tested}
18 discrete claims (C1--C18) enumerated in \texttt{REPORT.md}: 14 tested from public artifacts
(\checkmark), 2 phenotypic assays not attempted (no strain in hand; C15 cytotoxicity, C16
metal-tolerance MICs), 1 implicit (C17 CRISPR-Cas absence, inferred from PGAP), 1 cross-validation
(C18, via BV-BRC independent pipeline).

\section{Method}
Work performed on \textbf{uicgpu} (8$\times$A100, 255 cores, 2 TB RAM) at
\texttt{/data/stevens/BVBRC-92-PA34/}, with the LLM-judge step driven from CherryRd through the
Argo LLM proxy (\texttt{127.0.0.1:44497}).

\begin{enumerate}
  \item \textbf{Acquisition.} Paper PDF from PLoS (rasterized with \texttt{pdftotext -layout});
        all 6 genomes (CP032552, MH547560, MH547561, AE004091, CP000438, CP008739) via NCBI Entrez
        \texttt{efetch} (FASTA + GenBank). BV-BRC cross-reference via public REST API
        (\texttt{sp\_gene} table for \texttt{genome\_id} 287.6355).
  \item \textbf{Table 2 recomputation.} Direct Biopython 1.87 parse of each GenBank record:
        length + GC\% from \texttt{SeqIO.read(...).seq}; per-feature counts by iterating
        \texttt{rec.features}; protein FASTA emitted from \texttt{translation} qualifiers.
  \item \textbf{Pan-genome (Roary-style, DIAMOND+MCL).} Concatenate 23{,}555 proteins across
        4 genomes, all-vs-all BLASTP with DIAMOND 2.1.9 (\texttt{--more-sensitive}, e-val $\le$1e-5,
        246{,}824 hits), filter $\ge$50\% ID + $\ge$50\% coverage (78{,}801 hits), build undirected
        weighted graph (22{,}605 nodes / 39{,}655 edges), cluster with
        \texttt{markov\_clustering.run\_mcl} at inflation 1.5. Tag clusters by genome membership.
  \item \textbf{Per-locus AMR / virulence / mobilome verification.} Regex on
        \texttt{product}/\texttt{gene}/\texttt{note} qualifiers per CDS in each PA34 record;
        cross-check each hit's coordinate against the paper's Table 3 RGP intervals.
  \item \textbf{Independent BV-BRC cross-check.} Pull all \texttt{sp\_gene} records (Antibiotic
        Resistance + Metal Resistance) for \texttt{genome\_id} 287.6355; compare inventory.
  \item \textbf{LLM judge.} Assemble structured evidence bundle $\to$ \texttt{argo:gpt-5.2}
        (temperature 0.1) with an expert-bioinformatician grading prompt; capture JSON verdict.
\end{enumerate}

\section{Results vs paper}

\subsection{Table 2 --- chromosome statistics (CP032552)}
\begin{longtable}{@{}lrrl@{}}
\toprule
Field & Paper & Recomputed & Match \\
\midrule
Genome size (bp)       & 6{,}810{,}079 & 6{,}810{,}079 & EXACT \\
G+C \%                 & 66.1          & 66.07         & \checkmark \\
Genes                  & 6{,}544       & 6{,}544       & EXACT \\
CDS                    & 6{,}462       & 6{,}462       & EXACT \\
Protein-coding         & 6{,}314       & 6{,}314       & EXACT \\
Pseudogenes            & 148           & $\approx$148  & approx \\
Total RNA              & 82            & 81            & off-by-1 \\
tRNA                   & 65            & 65            & EXACT \\
ncRNA                  & 5             & 4             & off-by-1 \\
\bottomrule
\end{longtable}

\subsection{Plasmid statistics}
\begin{longtable}{@{}llrrl@{}}
\toprule
Plasmid & Field & Paper & Recomputed & Match \\
\midrule
pMKPA34-1 (MH547560) & Length     & 95.4 kbp & 95{,}404 bp & EXACT \\
                     & G+C \%     & 57       & 57.22       & \checkmark \\
                     & CDS        & 98       & 98          & EXACT \\
pMKPA34-2 (MH547561) & Length     & 26.8 kbp & 26{,}862 bp & EXACT \\
                     & G+C \%     & 61       & 61.00       & \checkmark \\
                     & CDS        & 33       & 32          & off-by-1 \\
\bottomrule
\end{longtable}

\subsection{Pan-genome (paper Roary vs our DIAMOND+MCL rerun)}
\begin{longtable}{@{}lrrrl@{}}
\toprule
Metric & Paper & Rerun & $\Delta$ & Match \\
\midrule
Pan (total orthologs)   & 7{,}643 & 6{,}775 & $-11.4\%$ & approx \\
Core (all 4)            & 5{,}078 & 4{,}654 & $-8.3\%$  & approx \\
\textbf{PA34 accessory} & \textbf{1{,}213} & \textbf{1{,}206} & \textbf{$-0.6\%$} & \textbf{ESSENTIALLY EXACT} \\
PA34 unique             & 543     & 661     & $+21.7\%$ & approx \\
PA34 no-ortho vs PAO1   & 886     & 855     & $-3.5\%$  & \checkmark \\
PA34 no-ortho vs PA14   & 737     & 701     & $-4.9\%$  & \checkmark \\
PA34 no-ortho vs VRFPA04& 946     & 1{,}007 & $+6.4\%$  & \checkmark \\
\bottomrule
\end{longtable}

The paper's directional finding---\textbf{VRFPA04 shares the fewest orthologs with PA34}, despite
also being ocular---reproduces. Absolute pan/core numbers diverge because our DIAMOND+MCL
clustering uses 50\% identity vs Roary's canonical 95\%, which inflates the cloud (singleton) count
and depresses core. Importantly, the paper's headline accessory-genome number reproduces to
within 1\%.

\subsection{Per-locus AMR / virulence / mobilome (chromosome CP032552)}
\begin{longtable}{@{}p{0.28\linewidth}p{0.30\linewidth}p{0.30\linewidth}l@{}}
\toprule
Feature & Paper claim & Our finding & \checkmark \\
\midrule
exoU (T3SS effector)          & RGP7 (4{,}719{,}909--4{,}727{,}427) & pos 4{,}720{,}713                  & \checkmark \\
SpcU chaperone                & (same island)                       & pos 4{,}720{,}303                  & \checkmark \\
\textbf{AAC(3)-IId}           & RGP23 (3{,}231{,}884--3{,}357{,}062)& pos 3{,}233{,}553                  & \checkmark \\
Tunicamycin resistance        & RGP23                               & pos 3{,}234{,}426                  & \checkmark \\
Copper resistance operon      & RGP23                               & 3{,}271{,}857--3{,}273{,}265 (copB) & \checkmark \\
Phage gp37                    & inserted into RGP23                 & pos 3{,}314{,}352                  & \checkmark \\
\textbf{Hg operon \#1 (MKPA34-GI1)} & 2{,}284{,}401--2{,}353{,}046 (novel GI) & mer proteins 2{,}342{,}693--2{,}345{,}779 & \checkmark \\
Chromate operon               & MKPA34-GI1                          & chrA at 2{,}298{,}255--2{,}299{,}135 & \checkmark \\
\textbf{Hg operon \#2 (RGP5)} & 5{,}010{,}479--5{,}090{,}313        & full merR-T-P-A-B-D at 5{,}075{,}000--5{,}080{,}589 & \checkmark \\
Pyoverdine (pvdE)             & RGP73 (3{,}022{,}522--3{,}055{,}524)& NRPS 2{,}988{,}776; export 3{,}061{,}910--3{,}065{,}331 & \checkmark (flanking) \\
Flagellin replacement         & RGP9 (4{,}605{,}383--4{,}612{,}609) & annotation at 4{,}600{,}486         & \checkmark (adjacent) \\
\bottomrule
\end{longtable}

\subsection{Plasmid AMR / mobilome}
\begin{itemize}
  \item \textbf{pMKPA34-1:} dfrA15 \checkmark, cmlA1 \checkmark, strA/APH(3\textquotedbl)-Ib
        \checkmark, strB/APH(6)-Id \checkmark, bla$_{\text{NPS-1}}$ \checkmark, sul1 \checkmark
        (bonus, canonical In1427 signature), intI1 \checkmark, Tn3 tnpA$\times$2 + tnpR \checkmark,
        acrA + acrB + oprM \checkmark, qacE$\Delta$1 \checkmark (bonus), repE/smc/parB/traN
        \checkmark, xerC/xerD \checkmark.
  \item \textbf{pMKPA34-2:} mepA \checkmark, tnsA/B/C/D/E \checkmark, phage integrase + resolvase
        \checkmark.
\end{itemize}

\subsection{Independent BV-BRC cross-check (\texttt{genome\_id} 287.6355)}
251 Antibiotic Resistance + 37 Metal Resistance annotations from BV-BRC's independent PATRIC
pipeline. Antibiotic hits include AAC(3)-II family, APH(6)-Id, CmlA, folA/Dfr, ampC/PDC-3, and
multiple \textit{mex} efflux operons. Metal hits include \textbf{merA$\times$2, merB$\times$2,
merP$\times$2, merR$\times$3}---completely independent evidence for two mercury-resistance operons,
matching the paper's central mobilome claim.

\subsection{LLM judge}
\texttt{argo:gpt-5.2}, temperature 0.1 --- verdict \textbf{PARTIAL} at \textbf{high} confidence.
One-line: \textit{``Genome/plasmid statistics and the major AMR/virulence/dual-mercury-operon
findings reproduce from public data, but the published pan-genome/core/unique gene counts are not
fully reproduced under an equivalent pipeline, so replication is partial.''}

\section{Verdict and justification}
\textbf{PARTIAL (strong).}

\paragraph{Fully replicated.} Chromosome length + GC + CDS + gene + tRNA + rRNA counts (Table 2);
plasmid pMKPA34-1 length + GC + CDS + full AMR gene inventory; plasmid pMKPA34-2 length + GC +
mepA + Tn7; PA34 accessory-genome headline (1{,}206 vs 1{,}213, $\Delta<1\%$); directional
finding that VRFPA04 shares fewest orthologs with PA34; all specific chromosomal AMR/virulence/metal
claims (exoU/SpcU in RGP7, AAC(3)-IId + tunicamycin + copper + phage gp37 in RGP23, two
independent mercury operons); BV-BRC pipeline agreement on the resistance inventory and on the
duplicated mer operon.

\paragraph{Partially replicated.} Absolute pan-genome / core-genome / PA34-unique counts
($-8$ to $+22\%$ vs paper) --- toolchain-parameter difference (DIAMOND+MCL at 50\% ID vs Roary
canonical 95\% ID). Not a biology difference. A faithful rerun would install Roary itself; time
budget did not permit here.

\paragraph{Not tested.} Fig 5 cytotoxicity (needs live strain + HCEC culture); Fig 6 MIC vs
Hg/Cu/Co (needs strain + wet lab); ST1284 MLST call (would need a dedicated PubMLST BLAST script
since PGAP annotations don't carry the traditional MLST gene tags).

\section{Genuine critique}
This section is the honest ``what did and did not replicate, and why the reader should still care''
--- not a summary.

\subsection{What actually replicated and matters}
\begin{itemize}
  \item \textbf{The deposited artifacts match the paper's headline statistics on the nose.} Length,
        GC, gene/CDS/tRNA counts for the chromosome are EXACT; both plasmids are exact on length
        and GC and off by at most one CDS. That is the strongest possible internal-consistency
        check: the sequence in the database is the sequence the paper describes.
  \item \textbf{Every named AMR / virulence / metal / mobilome locus is present at the stated
        coordinate.} Not ``elsewhere in the genome''---the actual RGP intervals contain the actual
        genes, including the two independent mercury operons in two different chromosomal regions.
        This is the paper's central biological claim and it survives independent verification
        with the coordinates preserved.
  \item \textbf{A completely independent annotation pipeline (BV-BRC's PATRIC) with a different
        assembly of the same isolate (128 SPAdes contigs vs the paper's closed hybrid assembly)
        reaches the same qualitative inventory}, including the mer operon duplication. Two
        pipelines, one signal.
\end{itemize}

\subsection{What did not replicate and why}
\begin{itemize}
  \item \textbf{Absolute pan-genome counts diverged by 8--22\%.} The gap is fully explained by
        clustering parameters: DIAMOND+MCL at 50\% identity is looser than Roary's canonical 95\%,
        which will split more marginal orthologs into singletons and shrink the core. The
        \textit{direction} of the paper's central pan-genome claim (PA34 shares fewest orthologs
        with VRFPA04) reproduces. The \textit{magnitude} of core/pan/unique does not, and this is
        the single most important caveat: without rerunning with Roary itself at the paper's
        thresholds, ``7{,}643 pan / 5{,}078 core / 543 unique'' is not a validated number, only a
        plausible one.
  \item \textbf{PA34-unique inflated by $+22\%$}, which is the expected side-effect of the looser
        clustering (marginal orthologs get pushed to singleton). This is a known Roary-vs-MCL
        artifact and not a signal of a real biological difference.
  \item \textbf{ncRNA count is off by 1 (5 paper vs 4 recomputed); pMKPA34-2 CDS is off by 1 (33 vs
        32).} These are annotation-version drift (paper used one PGAP/Prokka version; the deposited
        GenBank has since been updated). Small, real, and honestly reported here.
\end{itemize}

\subsection{What was out of scope and remains untested}
The two phenotypic figures (Fig 5 exoU cytotoxicity in HCECs; Fig 6 metal-tolerance MICs) require
the live isolate and BSL-2 cell culture. Neither was attempted. A reader who wants to trust the
functional claims (``PA34's exoU is functional''; ``PA34 is more Hg-tolerant than PAO1'') still
depends on the paper's wet-lab evidence---this replication cannot substitute for it. That is a
real limit of a public-data replication and is called out here.

\subsection{Where the paper is stronger than a quick read suggests}
The paper's core claim is \textbf{not} the pan-genome count. It is that PA34's inflated accessory
genome is packed with mobile-element-borne AMR and metal-resistance loci, plausibly reflecting its
ocular / MDR clinical origin. That claim is fully supported: two independent mercury operons on
the chromosome, a 125 kbp phage-derived AMR island (RGP23), a class-I integron plus Tn3 plus Tn7
across the two plasmids. Both direct GenBank parsing and BV-BRC's PATRIC pipeline confirm the
inventory. The Subedi et al.\ team deposited exactly what they said they had, and the biology
holds up.

\subsection{Where a fair reader should push back}
\begin{itemize}
  \item ``24 RGPs'' is a number produced by MAUVE on a specific 4-genome comparison. We did not
        re-run MAUVE; we verified the coordinates of the specific RGPs the paper names. A stricter
        replication would re-run MAUVE with the same reference set.
  \item ``PA34 accessory = 1{,}213'' reproduces to within 1\% only because our accessory definition
        (any cluster missing at least one of the 4 genomes) happens to be robust to clustering
        threshold. That is a happy coincidence, not a validation of the number.
  \item CRISPR-Cas absence (C17) is inferred from ``no Cas gene in the PGAP annotation'' rather
        than from a positive-negative run of CRISPRCasFinder. A rigorous replication would run
        that tool.
\end{itemize}

\subsection{Bottom line}
The paper's structural, AMR, virulence, and metal-resistance claims about PA34 replicate cleanly
from public data using an independent pipeline. The comparative-genomics counts replicate in
direction but not in magnitude, and the phenotypic claims are out of scope. Verdict PARTIAL
reflects both the successes and the honest gaps.

\section{Files}
\texttt{report/REPORT.md}, \texttt{report/brief.md}, \texttt{report/attempt\_log.md},
\texttt{report/artifact\_harvest.md}, \texttt{report/evidence/summary\_verification.json},
\texttt{report/evidence/pangenome\_result.json},
\texttt{report/evidence/bvbrc\_spgene\_pa34.json},
\texttt{report/evidence/genomes\_downloaded.txt},
\texttt{report/evidence/llm\_judge\_verdict.json} + \texttt{.txt},
\texttt{work/paper.pdf} + \texttt{work/paper.txt}, \texttt{work/pangenome\_pa34.py}.

\vspace{1em}
\noindent\textit{Independent replication conducted 2026-07-04 as part of the X-100 replication
project (BVBRC set, index 92). All data open-access (paper CC-BY 4.0 PLOS; sequences NCBI public;
BV-BRC free public API). No private data, no paywalled sources.}

\end{document}
